# Definitions for newly added A‑terms

| Term | Definition |
|------|------------|
| ACL Tear | A tear of the anterior cruciate ligament in the knee, usually caused by sports injuries, leading to instability and pain. |
| ADHD | Attention‑deficit/hyperactivity disorder, a neurodevelopmental condition characterized by inattention, hyperactivity, and impulsivity. |
| ADHD in Women | Presentation of ADHD in women, often with more inattentive symptoms and higher rates of internalizing disorders. |
| AIDS‑Defining Illnesses | Infections and cancers that, when occurring in HIV‑positive individuals, define the progression to AIDS (e.g., pneumocystis pneumonia, Kaposi sarcoma). |
| AL Amyloidosis (Primary Amyloidosis) | Deposition of immunoglobulin light‑chain amyloid fibrils in organs, causing organ dysfunction, especially kidneys and heart. |
| ALD (Adrenoleukodystrophy) | X‑linked disorder of peroxisomal fatty‑acid metabolism leading to adrenal insufficiency and progressive demyelination. |
| ARVC | Arrhythmogenic right‑ventricular cardiomyopathy, a genetic disease causing fibro‑fatty replacement of RV myocardium and ventricular arrhythmias. |
| ATTR‑CM (Transthyretin Amyloid Cardiomyopathy) | Amyloid deposition of transthyretin protein in the heart, causing restrictive cardiomyopathy and heart failure. |
| Abdominal Epilepsy | Rare form of epilepsy with seizure focus in the abdomen, causing unexplained abdominal pain or sensations. |
| Abdominal Muscle Strain | Overstretching or tearing of the abdominal wall muscles, often due to lifting or sudden movement. |
| Aberrant Right Subclavian Artery | Congenital anomaly where the right subclavian artery arises distal to the left artery, potentially causing dysphagia. |
| Abrasion (Scrape) | Superficial skin injury where the epidermis is scraped off, usually from friction. |
| Absence Seizures | Brief staring episodes with sudden loss of awareness, characteristic of childhood absence epilepsy. |
| Acetabular Fractures | Breaks of the acetabulum (hip socket) usually from high‑energy trauma, requiring orthopedic fixation. |
| Acetaminophen Toxicity & Overdose | Liver injury caused by excessive acetaminophen ingestion, leading to hepatic necrosis. |
| Achalasia (Cardiospasm) | Failure of the lower esophageal sphincter to relax, causing dysphagia and chest pain. |
| Achenbach Syndrome | Acute painful bruising of the finger due to capillary hemorrhage, self‑limited. |
| Achilles Paratenonitis | Inflammation of the paratenon surrounding the Achilles tendon, causing pain and swelling. |
| Acid Reflux & GERD | Backflow of gastric acid into the esophagus, causing heartburn and possible esophagitis. |
| Acne Keloidalis Nuchae | Chronic inflammatory papules and plaques on the nape of the neck, often in men of African descent. |
| Acne Papules | Small, raised inflammatory lesions in acne, consisting of a plugged pore with inflammation. |
| Acne Scars | Permanent depressions or raised tissue after severe acne lesions heal. |
| Acrophobia (Fear of Heights) | Intense, irrational fear of heights that may cause avoidance of high places. |
| Acute Cutaneous Lupus | Skin manifestation of lupus erythematosus presenting with photosensitive rash or lesions. |
| Acute Disseminated Encephalomyelitis (ADEM) | Inflammatory demyelinating disease of the brain and spinal cord, often post‑infectious. |
| Acute Fatty Liver of Pregnancy | Rare liver disorder in late pregnancy with microvesicular steatosis, potentially fatal. |
| Acute Heart Failure | Sudden onset of cardiac pump failure leading to pulmonary congestion and systemic hypoperfusion. |
| Acute Hepatic Porphyria | Acute attacks of porphyria caused by defects in hepatic heme synthesis, presenting with abdominal pain and neuropsychiatric symptoms. |
| Acute Lymphoblastic Leukemia (ALL) | Rapidly proliferating malignancy of lymphoid precursors in the bone marrow, common in children. |
| Acute Monocytic Leukemia | Subtype of acute myeloid leukemia characterized by proliferation of monocytic blasts. |
| Acute Myeloid Leukemia (AML) | Clonal proliferation of myeloid precursors leading to marrow failure and systemic symptoms. |
| Acute Respiratory Distress Syndrome (ARDS) | Severe lung injury with diffuse alveolar damage causing hypoxemia and respiratory failure. |
| Acyanotic Heart Disease | Congenital heart defects that do not cause cyanosis, such as atrial septal defect. |
| Addison’s Disease | Primary adrenal insufficiency causing deficiency of cortisol and aldosterone. |
| Adenocarcinoma Cancers | Malignant tumors arising from glandular epithelial tissue, found in many organs. |
| Adenoid Cystic Carcinoma (ACC) | Slow‑growing malignant tumor of salivary glands, often with perineural invasion. |
| Adie Syndrome (Tonic Pupil) | Neuropathy causing a dilated pupil with slow light reaction and absent deep tendon reflexes. |
| Adnexal Mass (Tumors) | Abnormal growth arising from ovarian or fallopian tube structures, benign or malignant. |
| Adrenal Adenoma | Benign tumor of the adrenal cortex, may be hormonally active. |
| Adrenal Gland Disorders | Various conditions affecting adrenal hormone production, including hyperplasia and tumors. |
| Adult Congenital Heart Disease | Congenital cardiac anomalies persisting into adulthood requiring lifelong management. |
| Adult‑Onset Still's Disease (AOSD) | Systemic inflammatory disease with high fevers, rash, arthritis, and elevated ferritin. |
| Advanced Maternal Age | Pregnancy in women aged 35 years or older, associated with increased obstetric risks. |
| Adventitial Cystic Disease | Cyst formation within the adventitia of an artery, most commonly the popliteal artery, causing limb ischemia. |
| Aerophobia (Fear of Flying) | Excessive anxiety about air travel that may lead to avoidance of flights. |
| Agenesis of the Corpus Callosum (ACC) | Absence of the major brain commissure, causing neurodevelopmental deficits. |
| Ageusia (Loss of Taste) | Complete loss of gustatory function, often due to nerve injury or infection. |
| Aichmophobia | Irrational fear of sharp objects such as needles or knives. |
| Ailurophobia (Fear of Cats) | Persistent anxiety and avoidance of cats. |
| Alcohol‑Induced Cardiomyopathy | Dilated cardiomyopathy resulting from chronic excessive alcohol consumption. |
| Alcohol‑Induced Hepatitis | Acute inflammation of the liver caused by excessive alcohol intake. |
| Alcohol‑Related Dementia | Cognitive decline associated with prolonged heavy alcohol use. |
| Alcohol‑Related Neuropathy | Peripheral nerve damage due to chronic alcohol abuse, causing sensory loss. |
| Alektorophobia (Fear of Chickens or Hens) | Specific phobia of chickens and related birds. |
| Algophobia (Fear of Pain) | Extreme fear of experiencing pain, often leading to avoidance of medical procedures. |
| Alice in Wonderland Syndrome | Distortions of visual perception, such as micropsia or macropsia, often linked to migraine. |
| Allergic Asthma | Asthma triggered by allergic reactions to inhaled allergens. |
| Allergic Rhinitis (Hay Fever) | IgE‑mediated inflammation of nasal mucosa causing sneezing, itching, and congestion. |
| Allergic Shiners | Dark circles under the eyes caused by venous congestion in allergic rhinitis. |
| Allergies | Hypersensitivity reactions mediated by IgE antibodies to environmental antigens. |
| Allergies: Shellfish | IgE‑mediated reaction to shellfish proteins, leading to urticaria, angioedema, or anaphylaxis. |
| Alpers Disease | Progressive neurodegenerative disorder with liver failure and seizures, often mitochondrial. |
| Alpha‑gal Syndrome | Delayed allergic reaction to red meat after tick bite, due to IgE to galactose‑α‑1,3‑galactose. |
| Altered Mental Status (AMS) | A spectrum of cognitive impairment ranging from confusion to coma. |
| Altitude Sickness | Acute illness from rapid ascent to high altitude, with headache, nausea, and dizziness. |
| Alzheimer’s Disease | Progressive neurodegenerative disorder causing memory loss, cognitive decline, and behavioral changes. |
| Amaxophobia (Fear of Driving) | Persistent fear of operating a motor vehicle, leading to avoidance. |
| Amblyopia (Lazy Eye) | Visual impairment of one eye due to abnormal visual development, treatable with patching. |
| Amebiasis (Amoebic Dysentery) | Intestinal infection with *Entamoeba histolytica* causing dysentery and liver abscesses. |
| Amplified Musculoskeletal Pain Syndrome (AMPS) in Children | Chronic pain disorder in children with exaggerated pain perception and disability. |
| Ampullary Cancer | Malignancy arising from the ampulla of Vater, often causing jaundice. |
| Amyoplasia | Congenital joint contractures with muscle hypoplasia, part of arthrogryposis. |
| Amyotrophic Lateral Sclerosis (ALS) | Progressive motor neuron disease causing muscle weakness and respiratory failure. |
| Anal Dysplasia | Precancerous changes in the anal epithelium, often associated with HPV infection. |
| Anal Fissures | Linear tear in the anoderm causing pain and bleeding during defecation. |
| Anal Herpes | Herpes simplex virus infection of the anal region, presenting with painful vesicles. |
| Anal Itching (Pruritus Ani) | Chronic itching of the perianal skin, often due to irritation or infection. |
| Anal Stenosis | Narrowing of the anal canal leading to constipation and pain. |
| Anal Warts | Condyloma acuminata of the anal region caused by HPV, presenting as papillomatous growths. |
| Anal Yeast Infection | *Candida* infection of the anal skin causing itching and erythema. |
| Anaplastic Thyroid Cancer (ATC) | Aggressive thyroid carcinoma with poor prognosis, often presenting with a rapidly enlarging neck mass. |
| Androphobia | Fear of men or masculinity, leading to avoidance of male individuals. |
| Anemia During Pregnancy | Reduced hemoglobin concentration in pregnancy, increasing risk for maternal and fetal complications. |
| Anemia Rash | Rare cutaneous manifestation of hemolytic anemia, presenting with erythematous lesions. |
| Anemia in Newborns | Low hemoglobin in neonates, often due to premature birth or hemolysis. |
| Anemia of Chronic Disease | Normocytic anemia associated with chronic inflammation, infection, or malignancy. |
| Aneurysmal Bone Cyst (ABC) | Expansile, blood‑filled bone lesion that can cause pain and swelling. |
| Angina (Angina Pectoris) | Chest discomfort due to myocardial ischemia, precipitated by exertion. |
| Anginophobia (Fear of Chest Pain) | Excessive anxiety about experiencing angina, leading to health‑care avoidance. |
| Angiodysplasia of Gastrointestinal (GI) Tract | Dilated tortuous vessels in the GI mucosa causing occult bleeding. |
| Angioid Streaks | Breaks in Bruch’s membrane of the eye associated with systemic diseases like PXE. |
| Angioimmunoblastic T‑cell Lymphoma (AITL) | Aggressive peripheral T‑cell lymphoma with systemic symptoms and skin rash. |
| Angiokeratoma | Benign vascular skin lesion, often dark red to black papules. |
| Angiolipoma | Benign tumor composed of fat and vascular tissue, typically subcutaneous. |
| Angiomatoid Fibrous Histiocytoma | Rare soft‑tissue tumor with blood‑filled spaces, may mimic malignancy. |
| Angiomyolipoma of the Kidney | Benign renal tumor containing blood vessels, smooth muscle, and fat; may bleed. |
| Anhidrosis (Lack of Sweat) | Absence of sweating, leading to overheating and heat intolerance. |
| Animal Bites | Injuries caused by animal teeth, risk of infection and need for prophylactic antibiotics. |
| Aniridia (Absence of Iris) | Congenital absence of the iris, often associated with foveal hypoplasia and cataracts. |
| Aniseikonia | Perceptual difference in image size between the two eyes, causing diplopia. |
| Anismus | Dysfunctional contraction of the pelvic floor muscles during defecation. |
| Anomic Aphasia | Language disorder characterized by difficulty retrieving words, while fluency and comprehension remain intact. |
| Anosmia (Loss of Sense of Smell) | Absence of olfactory perception, which may be congenital or acquired. |
| Ant Bites | Insect bites causing localized swelling, often allergic reaction. |
| Anterior Placenta | Placental implantation over the internal cervical os, increasing risk of bleeding. |
| Anteverted Uterus | Normal forward tilt of the uterus toward the bladder. |
| Anthophobia (Fear of Flowers) | Irrational fear of flowers, leading to avoidance. |
| Anthropophobia ( Fear of People) | Social phobia characterized by extreme anxiety around others. |
| Antidepressant Discontinuation Syndrome | Withdrawal symptoms after abrupt cessation of antidepressants, such as dizziness and flu‑like symptoms. |
| Antisocial Personality Disorder (ASPD) | Pattern of disregard for the rights of others, deceitfulness, and impulsivity. |
| Anxiety in Children | Excessive worry and physiological arousal in pediatric patients, affecting functioning. |
| Aortopathy | Disease of the aorta, including aneurysm, dissection, or coarctation. |
| Aphallia | Congenital absence of the penis. |
| Aphonia (Loss of Voice) | Inability to produce vocal sounds, often due to laryngeal nerve injury. |
| Apnea of Prematurity | Pauses in breathing in preterm infants due to immature respiratory control. |
| Appendicitis in Children | Inflammation of the appendix presenting with abdominal pain, fever, and leukocytosis in pediatric patients. |
| Appendicolith | Calcified deposit within the lumen of the appendix, associated with increased risk of perforation. |
| Aquagenic Pruritus | Intense itching triggered by contact with water, common in polycythemia vera. |
| Aquaphobia (Fear of Water) | Irrational fear of water, leading to avoidance of swimming or baths. |
| Arachnophobia (Fear of Spiders) | Persistent fear of spiders causing avoidance behaviors. |
| Arginine Vasopressin Disorders (Formerly Known as Nephrogenic Diabetes Insipidus) | Impaired renal response to vasopressin causing polyuria and polydipsia. |
| Argyria | Blue‑gray discoloration of skin due to chronic silver exposure. |
| Arithmophobia (Fear of Numbers) | Anxiety triggered by numbers or calculations. |
| Armpit Yeast Infection | Candidal infection of the axillary skin causing erythema and itching. |
| Arterial Thoracic Outlet Syndrome | Compression of subclavian artery in the thoracic outlet, causing upper‑extremity ischemia. |
| Arthritis in Hands | Inflammatory or degenerative disease affecting hand joints, causing pain and stiffness. |
| Arthritis in Knee | Joint inflammation of the knee, common in osteoarthritis and rheumatoid arthritis. |
| Arthritis in Wrist | Wrist joint inflammation leading to pain, swelling, and reduced motion. |
| Articulation Disorder | Speech disorder affecting the production of sounds due to motor planning deficits. |
| Asherman’s Syndrome | Intrauterine adhesions causing menstrual abnormalities and infertility. |
| Asphyxiation | Lack of oxygen intake leading to hypoxia and potential loss of consciousness. |
| Asplenia | Absence of splenic tissue, increasing susceptibility to encapsulated bacterial infections. |
| Asteatotic Eczema (Eczema Craquelé or Xerotic Eczema) | Dry, scaly dermatitis typical in the elderly, often on shins and forearms. |
| Asthma & Pregnancy | Management of asthma during pregnancy to maintain maternal oxygenation and fetal health. |
| Astraphobia (Fear of Thunder and Lightning) | Excessive anxiety during thunderstorms, leading to avoidance of outdoor activities. |
| Astrovirus | Viral cause of gastroenteritis, especially in children, leading to watery diarrhea. |
| Ataxia‑Teleangiectasia | Autosomal recessive disorder with cerebellar ataxia, immunodeficiency, and telangiectasias. |
| Ataxic Cerebral Palsy | Non‑progressive motor disorder with impaired balance and coordination due to brain injury. |
| Ataxophobia (Fear of Untidiness or Disorder) | Anxiety about disorderliness, often linked to obsessive‑compulsive tendencies. |
| Atelophobia (Fear of Imperfection) | Excessive perfectionism causing avoidance of tasks due to fear of making mistakes. |
| Atherosclerosis of Aorta | Plaque buildup in the aortic wall, increasing risk of aneurysm and embolism. |
| Athetoid Cerebral Palsy | Dyskinetic CP with involuntary writhing movements. |
| Athlete’s Foot (Tinea Pedis) | Fungal infection of the interdigital foot skin, causing itching and scaling. |
| Athlete’s Heart | Physiological cardiac remodeling in response to regular intense exercise. |
| Atrial Fibrillation (AFib) | Irregular rapid heart rhythm originating from the atria, increasing stroke risk. |
| Atrial Fibrillation With RVR | AFib with a rapid ventricular response (>100 bpm) requiring rate control. |
| Atrial Septal Aneurysm | Outpouching of the interatrial septum, sometimes associated with embolic events. |
| Atrioventricular Nodal Reentrant Tachycardia (AVNRT) | Fast supraventricular tachycardia due to a re‑entry circuit within the AV node. |
| Atrophic Kidney | Shrinkage and loss of renal parenchyma, often from chronic disease. |
| Atrophic Rhinitis | Chronic nasal mucosal atrophy leading to crusting and foul odor. |
| Attention‑Deficit/Hyperactivity Disorder (ADHD) in Adults | Persistent inattentive and hyperactive symptoms causing functional impairment in adulthood. |
| Atychiphobia (Fear of Failure) | Overwhelming dread of failing, often leading to avoidance of challenges. |
| Atypical (Walking) Pneumonia | Milder form of community‑acquired pneumonia with dry cough and subtle radiographic changes. |
| Atypical Anorexia | Restrictive eating disorder without body‑image disturbance. |
| Atypical Depression | Depressive symptoms without prominent sadness, often with irritability and somatic complaints. |
| Atypical Genitalia (Formerly Known as Ambiguous Genitalia) | Developmental variation of external genitalia not clearly male or female. |
| Atypical Teratoid/Rhabdoid Tumor (AT/RT) | Aggressive pediatric brain tumor with rhabdoid features. |
| Auditory Neuropathy | Disorder of auditory nerve transmission causing poor speech perception despite normal audiogram. |
| Auditory Processing Disorder (APD) | Difficulty processing auditory information in the brain, affecting comprehension. |
| Auto‑Brewery Syndrome | Endogenous ethanol production by gut microbes leading to intoxication without alcohol intake. |
| Autoimmune Inner Ear Disease (AIED) | Rapid progressive sensorineural hearing loss caused by autoimmune mechanisms. |
| Autoimmune Lymphoproliferative Syndrome (ALPS) | Defective Fas‑mediated lymphocyte apoptosis causing chronic lymphadenopathy and autoimmunity. |
| Autoimmune Neutropenia | Auto‑antibody mediated destruction of neutrophils, leading to recurrent infections. |
| Autonomic Dysreflexia (AD) | Life‑threatening hypertensive emergency in spinal cord injury patients triggered by noxious stimuli. |
| Autophobia (Fear of Being Alone) | Persistent anxiety about solitude, leading to constant need for companionship. |
| Avascular Necrosis (AVN or Osteonecrosis) | Death of bone tissue due to loss of blood supply, often affecting the femoral head. |
| Avoidant/Restrictive Food Intake Disorder (ARFID) | Eating disorder characterized by food avoidance without body‑image concerns, leading to weight loss or nutritional deficiency. |
| Axenfeld‑Rieger Syndrome | Developmental eye disorder with anterior segment anomalies and systemic features. |
| Axillo‑Subclavian Vein Thrombosis | Clot formation in the axillary or subclavian vein, causing upper‑extremity swelling. |

*Definitions are concise and suitable for a medical glossary.*
